\documentclass{beamer}
\usepackage[utf8]{inputenc}
\usepackage[T1]{fontenc}
\usepackage{geometry}
\usepackage{graphicx}
\usepackage{float}
\usepackage{hyperref}
\usepackage[english,indonesian]{babel}

\title{Optimasi Pemrosesan Template LaTeX Menggunakan Pendekatan Berbasis AI}
\author{Academic Support Agent}
\date{\today}

\begin{document}
\maketitle

\begin{abstract}

\end{abstract}

\chapter{PENDAHULUAN}


\chapter{TINJAUAN PUSTAKA}


% Original content (commented out):
% % Type of the document
\documentclass{beamer}

% elementary packages:
\usepackage{graphicx}
\usepackage[latin1]{inputenc}
\usepackage[T1]{fontenc}
\usepackage[english]{babel}
\usepackage{listings}
\usepackage{xcolor}
\usepackage{eso-pic}
\usepackage{mathrsfs}
\usepackage{url}
\usepackage{amssymb}
\usepackage{amsmath}
\usepackage{multirow}
\usepackage{hyperref}
\usepackage{booktabs}
\usepackage{tikz}



% additional packages
\usepackage{bbm}

% packages supplied with ise-beamer:
\usepackage{cooltooltips}
\usepackage{colordef}
\usepackage{beamerdefs}
\usepackage{lvblisting}

% Mathematics
\usepackage{amssymb}
\usepackage{amsmath}
\usepackage{mathrsfs}
\usepackage{amsthm,amsfonts}
\usepackage{mathtools}
\usepackage{algorithmic}
\usepackage[linesnumbered,ruled]{algorithm2e}
\usepackage{float}


% Change the pictures here:
% logobig and logosmall are the internal names for the pictures: do not modify them. 
% Pictures must be supplied as JPEG, PNG or, to be preferred, PDF
\pgfdeclareimage[height=2cm]{logobig}{Figures/hulogo}
% Supply the correct logo for your class and change the file name to "logo". The logo will appear in the lower
% right corner:
\pgfdeclareimage[height=0.7cm]{logosmall}{Figures/hulogo}

% Title page outline:
% use this number to modify the scaling of the headline on title page
\renewcommand{\titlescale}{1.0}
% the title page has two columns, the following two values determine the percentage each one should get
\renewcommand{\titlescale}{1.0}
\renewcommand{\leftcol}{0.6}

% smaller font for selected slides
\newcommand\Fontvi{\fontsize{10}{7.2}\selectfont}
\newcommand\Fontsm{\fontsize{8}{7.2}\selectfont}


% Define the title. Don't forget to insert an abbreviation instead 
% of "title for footer". It will appear in the lower left corner:
\title[Title shown at each slide]{Title for title page}
% Define the authors:
\authora{Author 1} % a-c
\authorb{Author 2}
\authorc{Author 3}

% Define any internet addresses, if you want to display them on the title page:
\def\linka{http://lvb.wiwi.hu-berlin.de}
\def\linkb{www.case.hu-berlin.de}
\def\linkc{}
% Define the institute:
\institute{Ladislaus von Bortkiewicz Chair of Statistics \\
C.A.S.E. -- Center for Applied Statistics\\
    and Economics\\
Humboldt--Universit{\"a}t zu Berlin \\}

% Comment the following command, if you don't want, that the pdf file starts in full screen mode:
\hypersetup{pdfpagemode=FullScreen}

%%%%
% Main document
%%%%
\begin{document}
% Draw title page
\frame[plain]{%
\titlepage{}
}

% The titles of the different sections of you talk, can be included via the \section command. The title will be displayed in the upper left corner. To indicate a new section, repeat the \section command with, of course, another section title
%%%%%%%%%%%%%%%%%%%%%%%%%%%%%%%%%%%%%%%%%%%%%%%%%%%%%%%%%%%%%%%%%%%%%%%%%%%%%%%%%%%%%%%%%%%%%%%%%%%%%%%%%%%%%%%%%%%%%%%%
%%%%%%%%%%%%%%%%%%%%%%%%%%%%%%%%%%%%%%%%%%%%%%%%%%%%%%%%%%%%%%%%%%%%%%%%%%%%%%%%%%%%%%%%%%%%%%%%%%%%%%%%%%%%%%%%%%%%%%%%
\frame{
\frametitle{Outline}
\begin{enumerate}
\item Introduction 
\item Pre-processing Steps
\item Model Selection
\item Variable Importance and Dimensionality Reduction 
\item Results and Conclusion
\end{enumerate}
}

\section{Introduction}
%%%%%%%%%%%%%%%%%%%%%%%%%%%%%%%%%%%%%%%%%%%%%%%%%%%%%%%%%%%%%%%%%%%%%%%%%%%%%%%%%%%%%%%%%%%%%%%%%%%%%%%%%%%%%%%%%%%%%%%%

% (A numbering of the slides can be useful for corrections, especially if you are
% dealing with large tex-files)
%%%%%%%%%%%%%%%%%%%%%%%%%%%%%%%%%%%%%%%%%%%%%%%%%%%%%%%%%%%%%%%%%%%%%%%%%%%%%%%%%%%%%%%%%%%%%%%%%%%%%%%%%%%%%%%%%%%%%%%%
\frame{
\frametitle{Formal Problem Setting}
\begin{itemize}
\item \textit{training set}: inputs $X = (x_1,\dots,x_n) \in \mathbb{R}^{n \times d}$ and labels $Y = (y_1,\dots,y_n)  \in  \mathbb{R}^{n}$
\item  \textit{test set}: inputs $X' = (x'_1,\dots,x'_t) \in \mathbb{R}^{t \times d}$ without labels
\end{itemize}
\vspace{0.5cm}
Find a function 
\begin{align}
f: X\rightarrow Y 
\end{align}
s.t. the \textit{test set} labels are predicted as accurately as possible, i.e.
\begin{align}
f(X') \approx Y'
\end{align} 
}

\section{Pre-Processing}
 
\frame{
\frametitle{Outline}
\begin{enumerate}
\item Introduction \quad \checkmark
\item Pre-processing Steps 
\item Model Selection
\item Variable Importance and Dimensionality Reduction 
\item Results and Conclusion
\end{enumerate}
}

\frame{
\vspace{0.1cm}
Several transformations and cleaning steps needed before putting the data into an algorithm, e.g. 
\frametitle{Pre-processing}

\begin{figure}
	\begin{center}
	\includegraphics[scale=0.25]{Figures/DataPipeline-1.jpg}
	\caption{Workflow of Pre-Processing Steps}
	\label{fig:DataPipeline}
	\end{center}
\end{figure}
All transformation need to be preformed on the test set as well! 
}

\begin{frame}[fragile]
\begin{center}
\begin{lstlisting}[
    basicstyle=\tiny, %or \small or \footnotesize etc.
]
basic_preprocessing = function(X_com, y, scaler="gaussian") 
{
	source("replace_ratings.R")
	source("convert_categoricals.R")
	source("impute_data.R")
	source("encode_time_variables.R")
	source("impute_outliers.R")
	source("scale_data.R")
	source("delete_nearzero_variables.R")
    X_ratings = replace_ratings(X_com)
    X_imputed = naive_imputation(X_ratings)
    X_no_outlier = data.frame(lapply(X_imputed, iqr_outlier))
    X_time_encoded = include_quarter_dummies(X_no_outlier)
    X_scaled = scale_data(X_time_encoded, scale_method = scaler)
    X_encoded = data.frame(lapply(X_scaled, cat_to_dummy))
    X_com = delect_nz_variable(X_encoded)
    idx_train = c(1:length(y))
    train = cbind(X_com[idx_train, ]
    test = X_com[-idx_train, ]
    return(list(train = train, X_com = X_com, test = test))
}
\end{lstlisting}
\end{center}
\quantnet \href{https://github.com/koehnden/SPL16/tree/master/Quantnet/dataProcessing/}{dataProcessing}
\end{frame}

%%%%%%%%%%%%%%%%%%%%%%%%%%%%%%%%%%%%%%%%%%%%%%%%%%%%%%%%%%%%%%%%%%%%%%%%%%%%%%%%%%%%%%%%%%%%%%
\section{Model Selection}

\frame{
\frametitle{Outline}
\begin{enumerate}
\item Introduction \quad \checkmark
\item Pre-processing Steps \quad \checkmark
\item Model Selection
\item Variable Importance and Dimensionality Reduction 
\item Results and Conclusion
\end{enumerate}
}


%%%%%%%%%%%%%%%%%%%%%%%%%%%%%%%%%%%%%%%%%%%%%%%%%%%%%%%%%%%%%%%%%%%%%%%%%%%%%%%%%%%%%%%%%%%%%%%%
\begin{frame}[fragile]
\frametitle{Optimizing Hyper-parameters} 

\begin{algorithm}[H]
\algsetup{linenosize=\tiny}
\scriptsize
\BlankLine
\ForEach{i in 1:t}
    {
     Randomly split the data into k folds of the same size \\
    	\ForEach{j in 1:k}
    	{
    	Use $j$th fold as test set and the union of remaining folds as training set \\
        \ForEach{p in 1:grid}
        {
                Fit model on training set using parameter set $p$ \\
                Predict on test set and calculate RMSE 
        }
    	}%end inner for
        \ForEach{p in 1:grid}{
                Calculate average RMSE over the $t \times k$-runs 
        }
        choose $p$ with the lowest RMSE
    }%end oute and r for
\caption{t-time k-fold crossvalidation and gridSearch}
\label{alg:seq}
\end{algorithm}


\quantnet \href{https://github.com/koehnden/SPL16/tree/master/Quantnet/xgbTuning/}{xgbTuning}
\quantnet \href{https://github.com/koehnden/SPL16/tree/master/Quantnet/rfTuning/}{rfTuning}
\quantnet \href{https://github.com/koehnden/SPL16/tree/master/Quantnet/svmTuning}{svmTuning}
\end{frame}



%%%%%%%%%%%%%%%%%%%%%%%%%%%%%%%%%%%%%%%%%%%%%%%%%%%%%%%%%%%%%%%%%%%%%%%%%%%%%%%%%%%%%%%%%%%%%%%%
\frame{
\frametitle{Taking on the curse of Dimensionality}
Problem: 
\begin{itemize}
\item many variables (99 after pre-processing)
\item small training set ($n = 1460$) 
\item variables are correlated with each other
\end{itemize}
\vspace{0.1cm}
Our approaches:
\begin{itemize}
\item Variable selection through variable importance ranking
\item Extract a smaller set of variable using PCA
\end{itemize}
}

\section{Results and Conclusion}
\frame{
\frametitle{Outline}
\begin{enumerate}
\item Introduction \quad \checkmark
\item Pre-processing Steps \quad \checkmark
\item Model Selection \quad \checkmark
\item Variable Importance and Dimensionality Reduction \quad \checkmark 
\item Results and Conclusion
\end{enumerate}
}

%%%%%%%%%%%%%%%%%%%%%%%%%%%%%%%%%%%%%%%%%%%%%%%%%%%%%%%%%%%%%%%%%%%%%%%%%%%%%%%%%%%%%%%%%%%%%%%%
\frame{
\frametitle{Results}
\begin{itemize}
\item Gaussian SVR with all variable is the single best model
\item PCA did not work well 
\item Models perform best with the full set of variables as Figure \ref{fig:RFE} suggested 
\end{itemize}
\vspace{0.25cm}
\begin{table}
\begin{center}
\begin{tabular}{c|ccc} 
\hline\hline
Inputs 		  & Gaussian SVR    & Random Forest & GBM      \\ 
\hline 
All Variables & \textbf{0.1308} &  0.1484       & 0.1333   \\
Top 30 		  & 0.1323  	 	&  0.1515       & 0.1436    \\
PCA	   		  & 0.1607  	    &  0.1657       & 0.1657     \\
\hline\hline
\end{tabular}
\caption{RMSE of submitted predictions}
\end{center}
\end{table}
\hspace{7.2cm} \href{https://github.com/koehnden/SPL16/blob/master/finalModels.R}{Github: finalModels}
}

\frame{
\frametitle{Outline}
\begin{enumerate}
\item Introduction \quad \checkmark
\item Pre-processing Steps \quad \checkmark
\item Model Selection \quad \checkmark
\item Variable Importance and Dimensionality Reduction \quad \checkmark 
\item Results and Conclusion \quad \checkmark 
\end{enumerate}
}


%%%%%%%%%%%%%%%%%%%%%%%%%%%%%%%%%%%%%%%%%%%%%%%%%%%%%%%%%%%%%%%%%%%%%%%%%%%%%%%%%%%%%%%%%%%%%%%%
% Dedicated section for references
\section{References}
\frame{
\frametitle{References}
\begin{thebibliography}{aaaaaaaaaaaaaaaaa}
\Fontvi
\beamertemplatearticlebibitems
\bibitem{Breiman:2003}
Breiman, Leo
\newblock{\em "Random Forest." Machine learning, 45(1), 5-32, (1999)}
\newblock available on \href{http://machinelearning202.pbworks.com/w/file/fetch/60606349/breiman_randomforests.pdf}{http://machinelearning202.pbworks.com}
\bibitem{ChenGuestrin:2015}
Chen, Tianqi, and Carlos Guestrin
\newblock{\em "XGBoost: Reliable Large-scale Tree Boosting System", Proceedings of the 22nd International Conference on Knowledge Discovery and Data Mining
Pages 785-794 (2015)}
\newblock available on \href{http://learningsys.org/papers/LearningSys_2015_paper_32.pdf}{http://learningsys.org} 
\beamertemplatearticlebibitems
\bibitem{DeCock:2011}
De Cock, Dean
\newblock{\em "Ames, Iowa: Alternative to the Boston housing data as an end of semester regression project" Journal of Statistics Education 19.3 (2011)}
\newblock available on \href{https://ww2.amstat.org/publications/jse/v19n3/decock.pdf}{https://ww2.amstat.org}
\beamertemplatearticlebibitems
\end{thebibliography}
}
%%%%%%%%%%%%%%%%%%%%%%%%%%%%%%%%%%%%%%%%%%%%%%%%%%%%%%%%%%%%%%%%%%%%%%%%%%%%%%%%%%%%%%%%%%%%%%%%%%%%%%%%
\frame{
\frametitle{References}
\begin{thebibliography}{aaaaaaaaaaaaaaaaa}
\Fontvi
\beamertemplatearticlebibitems
\bibitem{Friedman:2003}
Friedman, Jerome H.
\newblock{\em "Greedy function approximation: a gradient boosting machine." Annals of statistics 1189-1232 (2001).}
\newblock available on \href{http://projecteuclid.org/download/pdf_1/euclid.aos/1013203451}{https://www.jstor.org/journal/annalsstatistics}
\bibitem{Kuhn:2015}
Kuhn, Max, and Kjell Johnson
\newblock{\em "Applied predictive modeling". New York: Springer (2013)}
\beamertemplatearticlebibitems
\bibitem{Vapnik:1997}
Vapnik, Vladimir, Steven E. Golowich, and Alex Smola
\newblock{\em "Support vector method for function approximation, regression estimation, and signal processing." Advances in neural information processing systems 281-287 (1997)}
\newblock available on \href{https://pdfs.semanticscholar.org/43ff/a2c1a06a76e58a333f2e7d0bd498b24365ca.pdf}{https://semanticscholar.org}
\beamertemplatearticlebibitems
\end{thebibliography}
}



%%%%%%%%%%%%%%%%%%%%%%%%%%%%%%%%%%%%%%%%%%%%%%%%%%%%%%%%%%%%%%%%%%%%%%%%%%%%%%%%%%%%%%%%%%%%%%%%
\end{document}


\end{document}